Ex:- $T_1 : \mathbb{R}^3 \to \mathbb{R}^2$

$T_2 : \mathbb{R}^3 \to \mathbb{R}^2$

$T_1((x,y,z)) = (x+y, y+z)$

$T_2((x,y,z)) = (x, y+z)$

\begin{align*}
(T_1 + T_2)(x,y,z) &= T_1(x,y,z) + T_2(x,y,z) \\
&= (x+y, y+z) + (x, y+z) \\
&= (2x+y, 2y+2z)
\end{align*}

* $T : U \to V$

$S : V \to W$

\begin{align*}
S \circ T &= ST(u) \\
&= S(v) \\
&= w
\end{align*}

\section*{Nilpotent operator:-}

A linear transformation $T$ on a v-s $V$ is said to be nilpotent if $T^n = 0$, where $n$ is a least positive integer. where $n$ is known as the nilpotency index of $T$.

Ex:- Differential operator

$D : P_n \to P_n$

$D^{n+1} = 0$

2) $T : \mathbb{R}^3 \to \mathbb{R}^3$

\begin{align*}
T(x,y,z) &= (0,x,y) \\
T^2(x,y,z) &= T(0,x,y) = (0,0,x) \\
T^3(x,y,z) &= T(0,0,x) = (0,0,0)
\end{align*}

$\therefore$ nilpotency index of $T = 3$
